
This section introduces the general concepts behind a streaming lakehouse -- message queues, stream processing engines, open table formats, and dimensional modeling -- independent of any one product. Chapter~\ref{ch:proposed-solution} then explains which concrete tool implements each concept in this project, and why that tool was chosen over its alternatives.

\section{Why Real-Time Processing Matters}
\label{sec:why-realtime}

A modern platform such as a music-streaming service does not produce data in daily batches: it produces a continuous stream of small events -- a song played, a page viewed, a login -- the instant a user acts. Streamify, the project this report builds, simulates exactly this kind of platform with Eventsim, which generates synthetic ``song played'', ``page view'', and ``authentication'' events for a fake music website at a steady, ongoing rate.

In a traditional batch-processing setup, these events would simply accumulate in a source system and be pulled into the warehouse on a fixed schedule -- once an hour, once a day. Between runs, every dashboard and every downstream decision is working from data that is, on average, half of that interval out of date; right before the next run, it is a full interval out of date. For a platform that wants to react to what users are doing now -- surfacing trending songs, flagging a spike in failed logins, feeding a recommendation engine -- that lag is the problem: the data is correct, just always late.

Real-time (streaming) processing closes this gap: instead of waiting for a scheduled run, each event is picked up and made queryable within seconds to minutes of happening, so dashboards and downstream models reflect what is happening now rather than what happened last interval. In Streamify, this means a play event lands in its Kafka topic within milliseconds of Eventsim emitting it, and is queryable as a Hudi bronze row one Spark Structured Streaming micro-batch later (Figure~\ref{fig:batch-streaming}) -- typically a few seconds -- rather than at the next hourly or daily load.

Doing this reliably, and not just quickly, is harder than it sounds once the platform has more than one producer and more than one consumer of the same data. Four concrete problems show up:

\begin{itemize}
    \item \textbf{Ordering.} Events from different partitions of the same topic can arrive and be processed out of order relative to each other.
    \item \textbf{Fault tolerance.} Any component in the chain -- broker, processing job, storage layer -- can crash mid-stream and must be able to resume without losing data.
    \item \textbf{Exactly-once delivery.} On a retry after a crash, the same event must not be silently duplicated or dropped.
    \item \textbf{Schema drift.} Upstream event fields change shape over time (added, renamed, or retyped), and downstream tables must tolerate that without breaking.
    \item \textbf{Read/write contention.} The storage layer must serve a constant stream of small writes and large analytical reads at the same time, without one starving the other.
\end{itemize}

These are the distributed-systems problems that the rest of this section's concepts exist to solve: a \emph{message queue} for ordered, fault-tolerant ingestion; a \emph{stream processing engine} for continuous, exactly-once processing; and an \emph{open table format} for a storage layer that supports both frequent streaming writes and fast analytical reads.

\section{Message Queues and Publish-Subscribe Messaging}

\begin{figure}[H]
    \centering
    \includegraphics[width=0.7\textwidth]{2/message_queue.png}
    \caption{Producers and consumers decoupled through a publish-subscribe message queue.}
    \label{fig:message-queue}
\end{figure}

Event streaming treats data as an unbounded sequence of immutable events (e.g. ``user played a song'') rather than as static rows in a table. A \emph{message queue} built for this pattern is typically a log-based, publish-subscribe (pub-sub) broker, shown in Figure~\ref{fig:message-queue}: producers append events to a durable, ordered log rather than pushing directly to a fixed receiver, and publish them under named \emph{topics}; one or more consumers subscribe to the topics they need and read the log independently, at their own pace. This decouples producers from consumers in both time and number -- a producer does not need a consumer to be running or fast enough to keep up, and any number of independent consumers can read the same retained events without the broker deleting them on delivery.

\section{Batch and Streaming Processing Engines}

\begin{figure}[H]
    \centering
    \includegraphics[width=0.7\textwidth]{2/batch_steaming_processing.png}
    \caption{Batch processing versus streaming (event-at-a-time and micro-batching).}
    \label{fig:batch-streaming}
\end{figure}

A \emph{batch processing engine} operates on a bounded, complete dataset that is already at rest: it reads the full input, runs the job once end to end, produces its output, and stops until the next scheduled run.

A \emph{stream processing engine} instead operates on an unbounded dataset that keeps arriving: it stays running indefinitely, continuously reading new events from a message queue, applying the same kind of transformation logic (filtering, aggregation, joins) to them, and writing results incrementally as they are produced, rather than waiting for a full batch to accumulate first. Figure~\ref{fig:batch-streaming} contrasts the two.

A stream processing engine can still process its unbounded input using either of two execution models:

\begin{itemize}
    \item \textbf{Event-at-a-time.} Each event is processed as soon as it arrives, giving the lowest possible per-event latency.
    \item \textbf{Micro-batching.} The engine groups whatever new data has arrived since the last run into a small, discrete batch job on a fixed trigger interval, and processes that batch with the same execution engine used for ordinary batch queries. This reuses a mature batch engine's query planning and fault tolerance, at the cost of end-to-end latency being bounded by the trigger interval rather than being instantaneous.
\end{itemize}

Either execution model still reads its input positions (offsets) from the upstream message queue, and both keep the engine running continuously -- the distinction from batch processing is not how the input is chopped up, but that the engine never stops waiting for the next event.

\section{Storage layer \& Open Table Formats}
\label{sec:core-theory}

\begin{figure}[H]
    \centering
    \includegraphics[width=0.7\textwidth]{2/storage_layer.png}
    \caption{Object storage versus the table abstraction built on top of it.}
    \label{fig:storage-layer}
\end{figure}

As Figure~\ref{fig:storage-layer} shows, a lakehouse's underlying storage is one of three kinds:

\begin{itemize}
    \item \textbf{Object storage.} Data is stored as whole objects in flat buckets, addressed by key rather than by a filesystem path (e.g. S3, MinIO). This is the layer a lakehouse is built on.
    \item \textbf{Block storage.} Data is stored in fixed-size blocks that an operating system assembles into a filesystem, as used by a traditional database's local disk.
    \item \textbf{File storage.} Data is stored hierarchically as files in a shared, path-addressed filesystem (e.g. HDFS, NFS).
\end{itemize}

A lakehouse is built on object storage because it is cheap, horizontally scalable, and durable, but plain object storage is not a database: it does not provide row-level insert/update/delete semantics, ACID transactions, or time-travel queries -- it only stores and returns whole objects. An \emph{open table format} adds those capabilities on top of plain files in object storage, so a lake of files behaves like a table rather than an unmanaged pile of Parquet.

\begin{figure}[H]
    \centering
    \includegraphics[width=0.7\textwidth]{2/overview_open_table_format.png}
    \caption{Overview of an open table format, contrasting copy-on-write and merge-on-read.}
    \label{fig:open-table-format}
\end{figure}

Open table formats typically support two write strategies that trade off write cost against read cost, as summarized in Figure~\ref{fig:open-table-format}:

\begin{itemize}
    \item \textbf{COPY\_ON\_WRITE (COW).} Every write rewrites the affected data files as new versioned files. Reads are fast (plain columnar files, no merge step), but writes are more expensive because entire files are rewritten on update.
    \item \textbf{MERGE\_ON\_READ (MOR).} Writes are appended to a compact row-based delta log next to the existing base files; the two are merged only at read time (or during a background compaction job). Writes are cheap and low-latency, at the cost of extra work on read (or a periodic compaction step) to materialize the merged view.
\end{itemize}

MOR suits high-frequency streaming appends, where write latency matters most; COW suits tables that are rebuilt in full on a schedule and then read many times, where read performance matters most.

\section{Medallion Layering}

\begin{figure}[H]
    \centering
    \includegraphics[width=0.7\textwidth]{2/medalion_architecture.png}
    \caption{Medallion architecture: successive bronze-to-gold refinement layers.}
    \label{fig:medallion}
\end{figure}

The medallion architecture, shown in Figure~\ref{fig:medallion}, organizes a lakehouse into three successive layers of increasing refinement: bronze, silver, and gold.

\begin{itemize}
    \item \textbf{Bronze} tables hold raw or lightly-cleaned data close to its source form, written as-is by the ingestion layer.
    \item \textbf{Silver} tables hold data that has been cleaned, deduplicated, and conformed to a consistent schema, but not yet aggregated or reshaped for a specific consumer.
    \item \textbf{Gold} tables hold data that has been modeled and aggregated into the shape consumers (BI tools, dashboards) actually query, typically built by a transformation layer reading from silver (or, when a project skips silver, directly from bronze).
\end{itemize}

This project uses all three layers: dbt reads bronze and builds the five dimension tables as silver, then reads bronze together with those silver dimensions to build the gold \texttt{fact\_streams} and \texttt{wide\_streams} tables (Section~\ref{sec:transformation-layer}).
